\chapter{รากฐานฟิสิกส์การตรวจวัดค่ากรด-ด่างดินและผลกระทบอุณหภูมิความชื้น}
\label{chap:ch01}

\section*{แผนบริหารการสอนประจำบทที่ 1}
\addcontentsline{toc}{section}{แผนบริหารการสอนประจำบทที่ 1}

\begin{tcolorbox}[
    enhanced, breakable,
    colback=white, colframe=rbruNavy,
    fonttitle=\bfseries\color{white},
    title={1. หัวข้อเนื้อหาประจำบท},
    arc=1.5mm, boxrule=0.8pt, leftrule=3.5mm
]
\begin{enumerate}
    \item ความสำคัญยิ่งยวดของค่าความเป็นกรด-ด่างของดิน (Soil pH) ต่อการละลายและการดูดซึมธาตุอาหารพืช
    \item เคมีไฟฟ้าของการตรวจวัดค่ากรด-ด่างดินและสมการเนิร์นสต์ (Nernst Equation)
    \item กลไกความคลาดเคลื่อนจากอุณหภูมิดินต่อความชันศักย์ไฟฟ้าและการแตกตัวของน้ำ
    \item ผลกระทบจากปริมาณความชื้นในดินและอิมพีแดนซ์รอยต่อของเหลวในสภาวะดินแห้ง
    \item สถาปัตยกรรมการชดเชยค่าแบบสองขั้นตอนด้วยปัญญาประดิษฐ์อิงฟิสิกส์ (Two-Stage Decoupling Architecture)
\end{enumerate}
\end{tcolorbox}

\begin{tcolorbox}[
    enhanced, breakable,
    colback=white, colframe=rbruNavy,
    fonttitle=\bfseries\color{white},
    title={2. วัตถุประสงค์เชิงพฤติกรรม},
    arc=1.5mm, boxrule=0.8pt, leftrule=3.5mm
]
เมื่อศึกษาบทเรียนนี้จบแล้ว ผู้เรียนสามารถ
\begin{enumerate}
    \item อธิบายกลไกทางเคมีและสรีรวิทยาของรากพืชที่ตอบสนองต่อค่าความเป็นกรด-ด่างของดินได้อย่างถูกต้อง
    \item อธิบายสมการเนิร์นสต์และผลกระทบของอุณหภูมิต่อความชันศักย์ไฟฟ้าเคมีได้อย่างแม่นยำ
    \item วิเคราะห์สาเหตุของความคลาดเคลื่อนจากความชื้นต่ำและอิมพีแดนซ์รอยต่อของเหลวได้อย่างเป็นระบบ
    \item อธิบายหลักการสถาปัตยกรรมการชดเชยค่าสองขั้นตอนแบบ Edge Two-Stage Decoupling บนสมาร์ตโฟนได้อย่างสมบูรณ์
\end{enumerate}
\end{tcolorbox}

\begin{tcolorbox}[
    enhanced, breakable,
    colback=white, colframe=rbruNavy,
    fonttitle=\bfseries\color{white},
    title={3. วิธีสอนและกิจกรรมการเรียนการสอน},
    arc=1.5mm, boxrule=0.8pt, leftrule=3.5mm
]
\begin{enumerate}
    \item การบรรยายเชิงปฏิสัมพันธ์ร่วมกับการสาธิตปรากฏการณ์เคมีไฟฟ้าในสารละลายดิน
    \item การอภิปรายกลุ่มย่อยเกี่ยวกับกรณีศึกษาปัญหาดินกรดรุนแรงในแปลงทุเรียนจังหวัดจันทบุรีและตราด
    \item การฝึกปฏิบัติการคำนวณการเปลี่ยนแปลงความชันเนิร์นสเตียนตามอุณหภูมิและการปรับแก้ค่า pH
\end{enumerate}
\end{tcolorbox}

\begin{tcolorbox}[
    enhanced, breakable,
    colback=white, colframe=rbruNavy,
    fonttitle=\bfseries\color{white},
    title={4. สื่อการเรียนการสอน},
    arc=1.5mm, boxrule=0.8pt, leftrule=3.5mm
]
\begin{enumerate}
    \item สไลด์บรรยายอิเล็กทรอนิกส์และเอกสารคำสอนวิชาฟิสิกส์เกษตรดิจิทัล
    \item ชุดหัววัด Soil Parameter Tester (SoilpH-HT) และตัวอย่างสารละลายบัฟเฟอร์มาตรฐาน
    \item แอปพลิเคชันสมาร์ตโฟน SoilpH-HT พร้อมระบบจำลองข้อมูลเสมือนจริง
\end{enumerate}
\end{tcolorbox}

\begin{tcolorbox}[
    enhanced, breakable,
    colback=white, colframe=rbruNavy,
    fonttitle=\bfseries\color{white},
    title={5. การวัดและประเมินผล},
    arc=1.5mm, boxrule=0.8pt, leftrule=3.5mm
]
\begin{enumerate}
    \item การประเมินความรู้ผ่านแบบฝึกหัดและคำถามทบทวนท้ายบท
    \item การประเมินทักษะการคำนวณการชดเชยค่าความคลาดเคลื่อนของกรด-ด่างดิน
    \item การสังเกตพฤติกรรมการมีส่วนร่วมและการซักถามในชั้นเรียน
\end{enumerate}
\end{tcolorbox}

\clearpage

% ==============================================================================
% ภาพประกอบเกริ่นนำประจำบทที่ 1 (Introductory Infographic)
% ==============================================================================
\begin{figure}[H]
\centering
\begin{tikzpicture}[
    box/.style={rectangle, rounded corners=2.5mm, draw=rbruNavy, line width=1.0pt, fill=white, drop shadow={shadow xshift=1pt, shadow yshift=-1.5pt, opacity=0.15}, align=center, inner sep=6pt},
    header/.style={rectangle, rounded corners=1.5mm, draw=rbruNavy, fill=rbruNavy, text=white, font=\bfseries\small, align=center, inner sep=4pt},
    arrow/.style={-{Stealth[scale=1.1]}, line width=1.2pt, draw=rbruBlue}
]
    % คอลัมน์ที่ 1: แหล่งกำเนิดสัญญาณดิบจากหัววัด
    \node[box, fill=cyan!8, minimum width=3.6cm] (col1) {
        \textbf{สัญญาณตรวจวัดดิบ}\\[3pt]
        \footnotesize
        $\bullet$ กรด-ด่างดิบ ($\text{pH}_{\text{raw}}$)\\
        $\bullet$ อุณหภูมิดิน ($T_{\text{soil}}$)\\
        $\bullet$ ความชื้นดิน ($M_{\text{soil}}$)\\
        $\bullet$ สภาพนำไฟฟ้า ($EC_{\text{raw}}$)
    };
    \node[header, above=2pt of col1] {1. แหล่งกำเนิดสัญญาณ};

    % คอลัมน์ที่ 2: ทฤษฎีฟิสิกส์และเคมีไฟฟ้า
    \node[box, fill=rbruGold!10, minimum width=4.2cm, right=1.0cm of col1] (col2) {
        \textbf{สมการกำกับทางฟิสิกส์}\\[3pt]
        \footnotesize
        $\bullet$ Nernst Thermal Slope ($S_N$)\\
        $\bullet$ Water Dissociation ($pK_w$)\\
        $\bullet$ Liquid Junction Impedance\\
        $\bullet$ Topp \& Arrhenius Drift
    };
    \node[header, fill=rbruGold!90!black, above=2pt of col2] {2. กลไกฟิสิกส์เชิงลึก};

    % คอลัมน์ที่ 3: สถาปัตยกรรม Edge AI บนสมาร์ตโฟน
    \node[box, fill=green!8, minimum width=3.8cm, right=1.0cm of col2] (col3) {
        \textbf{ระบบ SoilpH-HT AI}\\[3pt]
        \footnotesize
        $\bullet$ USB-C OTG Driver\\
        $\bullet$ Modbus RTU RS485\\
        $\bullet$ Two-Stage PINN Engine\\
        $\bullet$ Calibrated pH + HUD
    };
    \node[header, fill=green!60!black, above=2pt of col3] {3. ปัญญาประดิษฐ์ฝังตัว};

    % ลูกศรเชื่อมโยงกระบวนการ
    \draw[arrow] (col1) -- node[above, font=\scriptsize\bfseries\color{rbruNavy}] {ส่งข้อมูลดิบ} (col2);
    \draw[arrow] (col2) -- node[above, font=\scriptsize\bfseries\color{rbruNavy}] {ชดเชยค่า PINN} (col3);

    % กรอบล่างแสดงเป้าหมายปลายทาง
    \node[rectangle, rounded corners=2mm, draw=rbruNavy!60, fill=rbruNavy!5, line width=0.8pt, minimum width=13.5cm, minimum height=0.9cm, below=0.7cm of col2, align=center] (goal) {
        \small\textbf{เป้าหมายทางวิชาการ} ยกระดับความแม่นยำในการวิเคราะห์ค่ากรด-ด่างดิน ตัดสัญญาณรบกวนจากอุณหภูมิและความชื้นแบบเรียลไทม์ ไร้สารเคมี
    };
    \draw[arrow, dashed] (col2.south) -- (goal.north);
\end{tikzpicture}
\caption{ผังมโนทัศน์เกริ่นนำประจำบทที่ 1 การบูรณาการฟิสิกส์การตรวจวัดค่ากรด-ด่างดิน ทฤษฎีเคมีไฟฟ้า และปัญญาประดิษฐ์บนสมาร์ตโฟน}
\label{fig:ch01_concept_overview}
\end{figure}

\section{วัตถุประสงค์ของการวิจัยและกรอบแนวคิด}
\label{sec:research_objectives}

โครงการวิจัยและพัฒนาระบบเครื่องมือวัดและแอปพลิเคชัน \textbf{SoilpH-HT AI} ดำเนินการภายใต้ระเบียบวิธีวิจัยฟิสิกส์เกษตรดิจิทัล มหาวิทยาลัยราชภัฏรำไพพรรณี โดยกำหนดวัตถุประสงค์หลัก 3 ประการให้สอดคล้องกับปัญหาจริงในภาคเกษตรกรรม ดังนี้

\begin{tcolorbox}[
    enhanced, breakable,
    colback=rbruNavy!5, colframe=rbruNavy,
    fonttitle=\bfseries\color{white},
    title={วัตถุประสงค์การวิจัย (Research Objectives)},
    arc=2.0mm, boxrule=1.0pt, leftrule=5.0mm
]
\begin{enumerate}
    \item \textbf{เพื่อศึกษาและสร้างชุดข้อมูลความสัมพันธ์ระหว่างค่าศักย์ไฟฟ้าของเซนเซอร์วัดค่าความเป็นกรด-เบสของดิน (Soil pH) กับค่าอุณหภูมิ และค่า Soil pH มาตรฐาน สำหรับพัฒนาโมเดลปัญญาประดิษฐ์} โดยทำการวัดค่าศักย์ไฟฟ้าของหัววัด ($E_{\text{sensor}}$ ในหน่วยมิลลิโวลต์) ควบคู่กับอุณหภูมิสารละลาย ($T$) ในสารละลายบัฟเฟอร์มาตรฐานปฐมภูมิของ NIST (pH 4.01, 7.00, 10.01) ในช่วงอุณหภูมิ $15^\circ\text{C}$ ถึง $45^\circ\text{C}$ เพื่อสร้างชุดข้อมูลการเรียนรู้ (Ground Truth Dataset) ที่แม่นยำสูง
    \item \textbf{เพื่อพัฒนาโมเดลปัญญาประดิษฐ์ที่สามารถชดเชยความคลาดเคลื่อนในการวัดที่เกิดจากผลของอุณหภูมิ (Temperature Effect) และความไม่เป็นเชิงเส้น (Non-linearity) ของเซนเซอร์} โดยออกแบบสถาปัตยกรรม Physics-Informed Neural Network (PINN) ที่ทำการแยกส่วนชดเชย (Decoupling) อย่างชัดเจนระหว่างผลจากอุณหภูมิตามสมการเนิร์นสต์และการเลื่อนของ $pK_w$ ($\Delta\text{pH}_{\text{temp}}$) ออกจากพฤติกรรมไม่เป็นเชิงเส้นและการตอบสนองแบบ Sub-Nernstian ของหัววัด ($\Delta\text{pH}_{\text{nonlinear}}$)
    \item \textbf{เพื่อทดสอบประสิทธิภาพของชุดทดสอบค่าความเป็นกรด-ด่างของดินแบบพกพาภาคสนามที่มีการชดเชยความคลาดเคลื่อนโดยปัญญาประดิษฐ์ โดยเปรียบเทียบความแม่นยำกับเครื่องมือวัดมาตรฐาน และทดสอบจริงในภาคสนาม} ในแปลงปลูกทุเรียนเศรษฐกิจในเขตจังหวัดจันทบุรีและจังหวัดตราด พร้อมวิเคราะห์ตามเกณฑ์มาตรฐานสากล ISO/IEC 17025 และ EURACHEM Guide
\end{enumerate}
\end{tcolorbox}

\section{ความสำคัญยิ่งยวดของค่าความเป็นกรด-ด่างของดิน}

การบริหารจัดการดินในงานเกษตรกรรมแม่นยำ (Precision Agriculture) โดยเฉพาะในแปลงปลูกพืชเศรษฐกิจมูลค่าสูง เช่น ทุเรียนหมอนทอง (\textit{Durio zibethinus} Murr.) ในเขตจังหวัดจันทบุรี ตราด และระยอง มีความจำเป็นต้องอาศัยข้อมูลสภาพดินที่ถูกต้อง โดยมีค่าความเป็นกรด-ด่างของดิน (Soil pH) เป็นตัวแปรควบคุมหลัก (Master Variable) ที่กำหนดพลวัตของธาตุอาหารและการเจริญเติบโตของพืช

\subsection{สรีรวิทยาของรากพืชและการตอบสนองต่อค่ากรด-ด่าง}
ระบบรากของทุเรียนเป็นรากที่ต้องการออกซิเจนสูงและมีความไวต่อสภาพความเป็นกรดในดินเป็นอย่างยิ่ง ช่วงค่าความเป็นกรด-ด่างที่เหมาะสมที่สุดต่อการเจริญเติบโตและการดูดซึมธาตุอาหารของทุเรียนอยู่ในช่วง pH 5.5 ถึง 6.5 หากดินมีความเป็นกรดรุนแรง (pH ต่ำกว่า 5.0) โครงสร้างทางเคมีและสรีรวิทยาของรากพืชจะได้รับผลกระทบอย่างรุนแรงดังนี้
\begin{enumerate}
    \item \textbf{ความเป็นพิษของอะลูมิเนียมและแมงกานีส (Aluminum and Manganese Toxicity)} เมื่อค่า pH ลดลงต่ำกว่า 5.0 สารประกอบอะลูมิโนซิลิเกตในอนุภาคดินจะแตกตัวปลดปล่อยไอออนอะลูมิเนียมอิสระ ($\text{Al}^{3+}$) และแมงกานีส ($\text{Mn}^{2+}$) เข้าสู่สารละลายดิน ไอออนเหล่านี้จะเข้าไปยับยั้งการแบ่งเซลล์และการยืดตัวของปลายรากฝอย ส่งผลให้รากกุด รากดำ ไม่สามารถดูดซึมน้ำและแร่ธาตุได้ตามปกติ
    \item \textbf{การตรึงธาตุฟอสฟอรัส (Phosphorus Fixation)} ในสภาวะดินกรด ไอออนฟอสเฟต ($\text{H}_2\text{PO}_4^-$) จะเข้าทำปฏิกิริยากับไอออนเหล็ก ($\text{Fe}^{3+}$) และอะลูมิเนียม เกิดเป็นสารประกอบเชิงซ้อนที่ไม่ละลายน้ำ (Insoluble Precipitates) ทำให้พืชขาดฟอสฟอรัส แม้เกษตรกรจะใส่ปุ๋ยฟอสเฟตในปริมาณมากก็ตาม ส่งผลให้การพัฒนาตาดอกและการติดผลล้มเหลว
    \item \textbf{การชะล้างธาตุอาหารรองและจุลธาตุ (Cation Leaching)} ประจุบวกของไฮโดรเจนไอออน ($\text{H}^+$) ปริมาณมหาศาลจะเข้าไปแย่งชิงตำแหน่งดูดซับบนผิวคอลลอยด์ของดิน ขับไล่ไอออนแคลเซียม ($\text{Ca}^{2+}$) แมกนีเซียม ($\text{Mg}^{2+}$) และโพแทสเซียม ($\text{K}^+$) หลุดออกจากอนุภาคดินและถูกชะล้างลึกลงสู่ชั้นดินล่าง ส่งผลให้เกิดอาการขาดธาตุอาหารอย่างรุนแรง
    \item \textbf{การหยุดชะงักของจุลินทรีย์ที่เป็นประโยชน์} สภาวะดินกรดรุนแรงจะยับยั้งการทำงานของแบคทีเรียกลุ่มตรึงไนโตรเจนและแบคทีเรียย่อยสลายอินทรียวัตถุ แต่กลับส่งเสริมการเจริญเติบโตของเชื้อราก่อโรครากเน่าโคนเน่า (\textit{Phytophthora palmivora}) ซึ่งเป็นโรคระบาดร้ายแรงอันดับหนึ่งในสวนทุเรียน
\end{enumerate}

\section{เคมีไฟฟ้าของการตรวจวัดค่ากรด-ด่างในดิน}

ค่าความเป็นกรด-ด่างของดิน นิยามตามหลักเคมีฟิสิกส์ว่าเป็นลอการิทึมฐานสิบของส่วนกลับกัมมันตภาพของไฮโดรเจนไอออน (Hydrogen Ion Activity, $a_{\text{H}^+}$) ในสารละลายดิน ดังสมการ
\begin{equation}
\text{pH} = -\log_{10}(a_{\text{H}^+}) = -\log_{10}(\gamma_{\text{H}^+} \cdot [\text{H}^+])
\end{equation}
โดยที่ $\gamma_{\text{H}^+}$ คือ สัมประสิทธิ์กัมมันตภาพ (Activity Coefficient) และ $[\text{H}^+]$ คือ ความเข้มข้นเชิงโมลของไฮโดรเจนไอออน

\subsection{สมการเนิร์นสต์และการกำเนิดศักย์ไฟฟ้าเคมีของเซนเซอร์}
ในการตรวจวัดค่ากรด-ด่างด้วยหัววัดทางเคมีไฟฟ้า ศักย์ไฟฟ้ารวมที่ตรวจวัดได้ ($E_{\text{cell}}$ หรือ $E_{\text{sensor}}$) เกิดจากผลต่างศักย์ไฟฟ้าระหว่างขั้ววัดแก้ว (Glass Indicator Electrode) กับขั้วอ้างอิงมาตรฐาน (Reference Electrode $\text{Ag/AgCl}$) ซึ่งอธิบายได้ด้วยสมการเนิร์นสต์ (Nernst Equation) สัมพันธ์กับค่า pH อุณหภูมิ และจุดสมศักย์ (Isopotential Point ที่ $\text{pH}_{\text{iso}} = 7.00$) ดังนี้
\begin{equation}
E_{\text{sensor}}(T, \text{pH}) = - S_N(T) \cdot (\text{pH} - \text{pH}_{\text{iso}}) = - S_N(T) \cdot (\text{pH} - 7.00)
\label{eq:sensor_potential_mv}
\end{equation}
โดยที่
\begin{itemize}
    \item $E_{\text{sensor}}$ คือ ศักย์ไฟฟ้าของขั้ววัด (Electrode Potential, หน่วยมิลลิโวลต์ mV)
    \item $S_N(T)$ คือ ความชันเนิร์นสเตียน (Nernstian Slope, หน่วย mV/pH) ที่อุณหภูมิสัมบูรณ์ $T$ (เคลวิน)
    \item $\text{pH}_{\text{iso}}$ คือ จุดความต่างศักย์เป็นศูนย์ (Isopotential Point ซึ่งโดยทั่วไปกำหนดที่ $\text{pH} = 7.00$ โดยที่ $E = 0\text{ mV}$)
\end{itemize}

ความชันเนิร์นสเตียนตามอุณหพลศาสตร์ $S_N(T)$ คำนวณได้จาก
\begin{equation}
S_N(T) = \frac{2.302585 \cdot R \cdot T}{F} \approx 54.197 + 0.1984 \times T_{\text{soil}} \quad (\text{mV/pH})
\label{eq:nernst_slope_calc}
\end{equation}
โดยที่
\begin{itemize}
    \item $R$ คือ ค่าคงที่สากลของก๊าซ ($8.31446 \, \text{J}\cdot\text{mol}^{-1}\cdot\text{K}^{-1}$)
    \item $T$ คือ อุณหภูมิสัมบูรณ์ ($T = T_{\text{soil}} [^\circ\text{C}] + 273.15\text{ K}$)
    \item $F$ คือ ค่าคงที่ฟาราเดย์ ($96,485.33 \, \text{C}\cdot\text{mol}^{-1}$)
\end{itemize}

ที่อุณหภูมิมาตรฐาน $25^\circ\text{C}$ ($298.15\text{ K}$) ความชันเนิร์นสเตียนทางทฤษฎีมีค่าเท่ากับ $59.16 \, \text{mV/pH}$ แต่เมื่ออุณหภูมิแปลงจริงพุ่งขึ้นเป็น $40^\circ\text{C}$ ความชันจะเพิ่มขึ้นเป็น $62.14 \, \text{mV/pH}$ ดังนั้น สารละลายที่มี $\text{pH} = 4.00$ จะสร้างศักย์ไฟฟ้าเคมีเท่ากับ $+177.48\text{ mV}$ ที่ $25^\circ\text{C}$ และเปลี่ยนเป็น $+186.42\text{ mV}$ ที่ $40^\circ\text{C}$ เกิดผลต่างศักย์ถึง $+8.94\text{ mV}$ ซึ่งหากไม่มีการแปลงและชดเชยที่ถูกต้องจะทำให้ระบบอ่านค่าผิดพลาดไปมากกว่า $0.15 - 0.25\text{ pH}$ ทันที

\section{แหล่งกำเนิดความคลาดเคลื่อนทางฟิสิกส์จากอุณหภูมิและความชื้น}

ในการนำหัววัดไปตรวจวัดในสภาพแปลงเกษตรกรรมจริง ค่าศักย์ไฟฟ้าที่หัววัดอ่านได้มักเกิดความคลาดเคลื่อนอย่างมีนัยสำคัญอันเนื่องมาจากปัจจัยสิ่งแวดล้อม 2 ประการหลัก ได้แก่

\subsection{1. ความคลาดเคลื่อนจากการเลื่อนลอยตามอุณหภูมิ ($\Delta \text{pH}_T$)}
อุณหภูมิดินในแปลงเกษตรกรรมเขตร้อนผันผวนอย่างรุนแรงตามรอบวัน ตั้งแต่ $22^\circ\text{C}$ ในช่วงเช้าตรู่ จนถึงมากกว่า $42^\circ\text{C}$ ในช่วงบ่าย ความผันผวนของอุณหภูมิก่อให้เกิดความคลาดเคลื่อน 2 ส่วนพร้อมกัน
\begin{enumerate}
    \item \textbf{การเปลี่ยนแปลงความชันเนิร์นสเตียนตามอุณหภูมิ} เมื่ออุณหภูมิสูงขึ้นเป็น $40^\circ\text{C}$ ($313.15\text{ K}$) ความชันจะเพิ่มขึ้นเป็น $62.14 \, \text{mV/pH}$ หากไม่มีการชดเชยอุณหภูมิ การวัดสารละลายที่มี pH 4.0 หรือ pH 9.0 จะเกิดความคลาดเคลื่อนสูงถึง $0.2 - 0.4 \, \text{pH}$ ทันที
    \item \textbf{การเลื่อนของค่าคงที่การแตกตัวของน้ำ ($pK_w$)} น้ำในสารละลายดินมีการแตกตัวเป็นไอออนตามอุณหพลศาสตร์ โดยค่า $pK_w$ จะลดลงเมื่ออุณหภูมิสูงขึ้น ส่งผลให้จุดเป็นกลางทางเคมีเลื่อนออกจาก 7.00
\end{enumerate}

ความคลาดเคลื่อนจากอุณหภูมิรวมจำลองได้ด้วยสมการฟิสิกส์
\begin{equation}
\Delta \text{pH}_T = (\text{pH}_{\text{raw}} - 7.0)\left(\frac{298.15}{273.15 + T_{\text{soil}}} - 1.0\right) + 0.017 \times (25.0 - T_{\text{soil}})
\end{equation}

\subsection{2. ความคลาดเคลื่อนจากอิมพีแดนซ์รอยต่อในสภาวะดินแห้ง ($\Delta \text{pH}_M$)}
ในดินที่มีความชื้นเชิงปริมาตร ($M_{\text{soil}}$) ต่ำกว่าร้อยละ 25 อนุภาคดินจะขาดฟิล์มน้ำต่อเนื่อง ส่งผลให้ความต้านทานสัมผัสระหว่างผิวหัววัดกับสารละลายดินพุ่งสูงขึ้นอย่างรวดเร็ว ก่อให้เกิดอิมพีแดนซ์รอยต่อของเหลว (Liquid Junction Impedance) และเกิดศักย์ไฟฟ้าอสมมาตร (Asymmetric Potential) ทำให้ค่า pH ที่อ่านได้เบี่ยงเบนจากค่าจริงอย่างมาก แบบจำลองทางฟิสิกส์แสดงความสัมพันธ์เชิงเอกซ์โพเนนเชียล
\begin{equation}
\Delta \text{pH}_M = \alpha \cdot \exp(-\beta \cdot M_{\text{soil}})
\end{equation}
โดยที่ $\alpha = 0.35$ และ $\beta = 0.08$ เป็นค่าสัมประสิทธิ์เฉพาะของดินร่วนปนทรายและดินเหนียวในภาคตะวันออก

\begin{table}[H]
\centering
\caption{การเปรียบเทียบสมรรถนะระหว่างวิธีการตรวจวัดดินแบบดั้งเดิมกับระบบ SoilpH-HT AI}
\label{tab:comparison_methods}
\begin{tabularx}{\textwidth}{p{3.0cm}p{3.2cm}p{3.5cm}Y}
    \toprule
    \textbf{คุณลักษณะ} & \textbf{วิธีแล็บเคมีมาตรฐาน} & \textbf{ชุดทดสอบเคมีภาคสนาม} & \textbf{ระบบ SoilpH-HT AI} \\
    \midrule
    ระยะเวลาวิเคราะห์ & 7 ถึง 14 วัน & 15 ถึง 30 นาที & เรียลไทม์ ($< 0.2$ วินาที) \\
    ค่าใช้จ่ายต่อตัวอย่าง & 500 ถึง 1,500 บาท & 50 ถึง 100 บาท & ไม่มีค่าใช้จ่ายซ้ำซ้อน \\
    ความแม่นยำเชิงปริมาณ & สูงมาก ($> 98\%$) & ปานกลางถึงต่ำ ($< 70\%$) & สูงมาก ($> 96\%$ ด้วย PINN) \\
    การชดเชยอุณหภูมิ & มีระบบปรับสภาพห้อง & ไม่มี & อัตโนมัติด้วย Deep Learning \\
    การชดเชยความชื้นดิน & ตากดินแห้งก่อนวัด & ไม่มี & ชดเชยแบบเรียลไทม์ในแปลง \\
    ผลกระทบสิ่งแวดล้อม & มีของเสียสารเคมี & เกิดขยะสารเคมีอันตราย & ไร้สารเคมีและเป็นมิตร 100\% \\
    \bottomrule
\end{tabularx}
\end{table}

\section{เกณฑ์มาตรฐานแถบสีค่ากรด-ด่างดินและการแปลผลทางปฐพีวิทยา}

เพื่อเชื่อมโยงผลการวัดเชิงตัวเลขของระบบดิจิทัลเข้ากับความรู้ความเข้าใจของเกษตรกรและนักวิชาการเกษตร ระบบ \textbf{SoilpH-HT} ได้ออกแบบตัวแปลงสีดิจิทัล (Digital Colorimetric Bridge) ตามเกณฑ์มาตรฐานแถบสีของกรมพัฒนาที่ดิน (LDD) และองค์การอาหารและการเกษตรแห่งสหประชาชาติ (FAO) ดังแสดงในตารางที่ \ref{tab:soil_ph_color_scale}

\begin{table}[H]
\centering
\caption{เกณฑ์มาตรฐานแถบสีค่าความเป็นกรด-ด่างของดิน (Soil pH) และการแปลผลทางสรีรวิทยาพืช}
\label{tab:soil_ph_color_scale}
\begin{tabularx}{\textwidth}{p{2.6cm}p{1.8cm}p{2.8cm}p{1.8cm}Y}
    \toprule
    \textbf{ระดับ pH} & \textbf{ช่วงค่า} & \textbf{แถบสีสารละลาย} & \textbf{รหัสสี HEX} & \textbf{ผลกระทบทางธาตุอาหารพืช} \\
    \midrule
    กรดรุนแรงมาก & $< 4.0$ & แดงเข้มจัด & \texttt{\#D32F2F} & อะลูมิเนียมและแมงกานีสเป็นพิษ รากพืชชะงักงัน \\
    กรดจัดมาก & $4.0 - 4.5$ & แดงส้มสด & \texttt{\#F4511E} & ฟอสฟอรัสถูกตรึงสูงมาก จุลินทรีย์ดินหยุดทำงาน \\
    กรดจัด & $4.5 - 5.5$ & ส้มเหลือง & \texttt{\#FB8C00} & ธาตุหลัก N, P, K ถูกดูดซึมได้ต่ำ ต้องใส่ปูนโดโลไมต์ \\
    กรดเล็กน้อย & $5.5 - 6.0$ & เหลืองอมเขียว & \texttt{\#FDD835} & พืชเขตร้อนส่วนใหญ่เจริญเติบโตได้ปานกลาง \\
    \textbf{กรดอ่อนๆ} & $\mathbf{6.0 - 6.5}$ & \textbf{เขียวตองอ่อน} & \texttt{\#C0CA33} & \textbf{จุดสมบูรณ์แบบ ทุเรียนดูดซึมปุ๋ยได้สูงสุด} \\
    \textbf{เป็นกลาง} & $\mathbf{6.5 - 7.5}$ & \textbf{เขียวมรกตสด} & \texttt{\#43A047} & \textbf{ระดับสมดุล ธาตุอาหารละลายตัวดีที่สุด} \\
    ด่างปานกลาง & $7.5 - 8.5$ & ฟ้าคราม & \texttt{\#0288D1} & ฟอสฟอรัสเริ่มตกตะกอนร่วมกับแคลเซียม \\
    ด่างจัดมาก & $> 8.5$ & ม่วงครามเข้ม & \texttt{\#5E35B1} & ดินเค็มโซดิก ดินแน่นทึบ ขาดธาตุเหล็กและสังกะสี \\
    \bottomrule
\end{tabularx}
\end{table}

\section{เกณฑ์การคำนวณอัตราปูนโดโลไมต์ปรับปรุงดินตามเนื้อดิน}

เมื่อตรวจพบค่าความเป็นกรดในดินต่ำกว่า 5.5 ระบบจะคำนวณใบสั่งอัตราปูนโดโลไมต์ ($\text{CaMg}(\text{CO}_3)_2$) หรือปูนขาวปรับปรุงดินโดยอัตโนมัติ โดยจำแนกตามเนื้อดิน (Soil Texture) ดังแสดงในตารางที่ \ref{tab:lime_requirement}

\begin{table}[H]
\centering
\caption{อัตราการใส่ปูนโดโลไมต์เพื่อยกระดับค่าความเป็นกรด-ด่างของดินขึ้น 1.0 หน่วย (กิโลกรัมต่อไร่)}
\label{tab:lime_requirement}
\begin{tabularx}{\textwidth}{p{3.2cm}cccY}
    \toprule
    \textbf{ระดับความเป็นกรดปัจจุบัน} & \textbf{ดินทราย / ร่วนปนทราย} & \textbf{ดินร่วน / ดินร่วนเหนียว} & \textbf{ดินเหนียวจัด} & \textbf{คำแนะนำทางปฐพีวิทยา} \\
    \midrule
    pH ต่ำกว่า 4.5 (กรดจัดมาก) & 250 ถึง 350 & 400 ถึง 550 & 600 ถึง 800 & แบ่งใส่ 2 ครั้ง ห่างกัน 1 เดือน \\
    pH 4.5 ถึง 5.0 (กรดจัด) & 180 ถึง 250 & 300 ถึง 400 & 450 ถึง 600 & หว่านรอบทรงพุ่มแล้วให้น้ำตาม \\
    pH 5.0 ถึง 5.5 (กรดปานกลาง) & 100 ถึง 180 & 200 ถึง 300 & 300 ถึง 450 & ใส่บำรุงประจำปีช่วงต้นฤดูฝน \\
    pH 5.5 ถึง 6.5 (เหมาะสม) & 0 & 0 & 0 & ดินสมบูรณ์ ไม่จำเป็นต้องใส่ปูน \\
    \bottomrule
\end{tabularx}
\end{table}

\section{คุณสมบัติของดินและสภาพแวดล้อมที่เหมาะสมสำหรับกลุ่มพืชเศรษฐกิจ}

การตรวจวัดและวินิจฉัยคุณภาพดินอย่างแม่นยำ จำเป็นต้องอาศัยข้อมูลเกณฑ์อ้างอิงคุณสมบัติทางกายภาพและเคมีที่เหมาะสมกับชีวสรีรวิทยาของพืชแต่ละกลุ่ม เนื่องจากพืชแต่ละชนิดมีระดับความทนทานและความต้องการธาตุอาหาร อุณหภูมิ และความชื้นที่แตกต่างกันอย่างมีนัยสำคัญ

\subsection{ปัจจัยสภาพแวดล้อมทางปฐพีวิทยาและบรรยากาศ}

ในการทำการเกษตรแม่นยำ ปัจจัยสภาพแวดล้อมที่มีอิทธิพลโดยตรงต่อการเจริญเติบโต การสร้างผลผลิต และการดูดซึมธาตุอาหาร ประกอบด้วย
\begin{enumerate}
    \item \textbf{ค่าความเป็นกรด-ด่างของดิน (Soil pH)} เป็นตัวแปรควบคุมหลัก (Master Variable) ที่กำหนดสภาพการละลายได้ของสารประกอบไอออนิกในสารละลายดิน โดยช่วง pH 6.0 ถึง 6.5 ถือเป็นจุดสมดุลที่ดีที่สุดสำหรับพืชเขตร้อนส่วนใหญ่
    \item \textbf{อุณหภูมิดินและอุณหภูมิอากาศ ($T_{\text{soil}}$ และ $T_{\text{air}}$)} อุณหภูมิดินที่เหมาะสมอยู่ในช่วง 24 ถึง 32 องศาเซลเซียส ส่งผลโดยตรงต่ออัตราการหายใจของรากพืช การแบ่งเซลล์เนื้อเยื่อเจริญ และกิจกรรมของจุลินทรีย์ดิน หากอุณหภูมิดินสูงเกิน 35 องศาเซลเซียส จะส่งผลให้รากพืชชะงักการทำงานและเสี่ยงต่อการเกิดความเครียดจากความร้อน (Heat Stress)
    \item \textbf{ความชื้นในดินและความชื้นสัมพัทธ์ในอากาศ (Soil Moisture และ Air RH)} ปริมาณน้ำในดินเชิงปริมาตร (VWC) ที่เหมาะสมอยู่ระหว่างร้อยละ 50 ถึง 70 สำหรับพืชบกทั่วไป ซึ่งเป็นช่วงความชื้นที่อนุภาคดินมีน้ำพร้อมใช้ (Available Water) โดยไม่เกิดสภาพขาดออกซิเจน ขณะที่ความชื้นสัมพัทธ์ในอากาศร้อยละ 60 ถึง 80 ช่วยรักษาสมดุลแรงดันไอระเหย (Vapor Pressure Deficit: VPD) ป้องกันการคายน้ำมากเกินควร
    \item \textbf{ธาตุอาหารหลัก (Primary Macronutrients: N-P-K)} ไนโตรเจน (N) ส่งเสริมการเจริญเติบโตทางกิ่งก้านและใบ ฟอสฟอรัส (P) พัฒนาระบบรากและการสร้างตาดอก และโพแทสเซียม (K) ควบคุมการเปิดปิดปากใบ การสังเคราะห์แป้งและน้ำตาลในผลผลิต
    \item \textbf{ธาตุอาหารรอง (Secondary Macronutrients: Ca-Mg-S)} แคลเซียม (Ca) สร้างความแข็งแรงแก่โครงสร้างผนังเซลล์และเยื่อหุ้มเซลล์ แมกนีเซียม (Mg) เป็นอะตอมแกนกลางของคลอโรฟิลล์ในการสังเคราะห์ด้วยแสง และกำมะถัน (S) เป็นองค์ประกอบของกรดอะมิโนจำเป็นและสารให้กลิ่นหอมเฉพาะตัว
    \item \textbf{จุลธาตุอาหารเสริม (Micronutrients: Fe-Mn-Zn-Cu-B-Mo)} แม้พืชต้องการในปริมาณระดับมิลลิกรัมต่อกิโลกรัม แต่ขาดไม่ได้ โดยโบรอน (B) ควบคุมการงอกของละอองเกสรและการติดผล สังกะสี (Zn) ควบคุมการสร้างฮอร์โมนออกซิน และเหล็ก (Fe) ร่วมในกระบวนการถ่ายทอดอิเล็กตรอน
\end{enumerate}

\subsection{เกณฑ์มาตรฐานคุณสมบัติดินและสภาพแวดล้อมจำแนกตามชนิดพืชเศรษฐกิจ}

เพื่อใช้เป็นฐานข้อมูลอ้างอิงในระบบ AI และระบบออกใบสั่งจัดการแปลง ตารางที่ \ref{tab:crop_soil_physical_npk} แสดงคุณสมบัติทางฟิสิกส์ สภาพบรรยากาศ และธาตุอาหารหลัก และตารางที่ \ref{tab:crop_soil_secondary_micro} แสดงธาตุอาหารรองและจุลธาตุสำหรับกลุ่มนาข้าว พืชไร่ และไม้ผลเศรษฐกิจ

\begin{table}[H]
\centering
\caption{เกณฑ์คุณสมบัติทางฟิสิกส์ สภาพแวดล้อม และธาตุอาหารหลัก NPK ที่เหมาะสมสำหรับกลุ่มพืชเศรษฐกิจ}
\label{tab:crop_soil_physical_npk}
\resizebox{\textwidth}{!}{%
\begin{tabularx}{1.18\textwidth}{p{2.8cm}ccccYYY}
    \toprule
    \textbf{กลุ่มและชนิดพืช} & \textbf{ช่วง pH ที่เหมาะสม} & \textbf{อุณหภูมิดิน ($^\circ\text{C}$)} & \textbf{อุณหภูมิอากาศ ($^\circ\text{C}$)} & \textbf{ความชื้นดิน (\%VWC)} & \textbf{ไนโตรเจน N (mg/kg)} & \textbf{ฟอสฟอรัส P (mg/kg)} & \textbf{โพแทสเซียม K (mg/kg)} \\
    \midrule
    \multicolumn{8}{l}{\textbf{1. กลุ่มนาข้าว (Rice Ecosystem)}} \\
    ข้าวหอมมะลิ 105 & $5.5 - 6.5$ & 25 ถึง 32 & 24 ถึง 33 & 70 ถึง 100 (ขังน้ำ) & 15 ถึง 30 & 10 ถึง 25 & 60 ถึง 100 \\
    ข้าวเจ้าชลประทาน & $5.5 - 7.0$ & 26 ถึง 32 & 25 ถึง 35 & 80 ถึง 100 (ขังน้ำ) & 20 ถึง 40 & 15 ถึง 30 & 70 ถึง 120 \\
    \midrule
    \multicolumn{8}{l}{\textbf{2. กลุ่มพืชไร่ (Field Crops)}} \\
    ข้าวโพดเลี้ยงสัตว์ & $5.8 - 6.8$ & 22 ถึง 30 & 20 ถึง 32 & 50 ถึง 70 & 30 ถึง 50 & 20 ถึง 40 & 80 ถึง 140 \\
    มันสำปะหลัง & $5.0 - 6.5$ & 25 ถึง 35 & 25 ถึง 35 & 30 ถึง 50 & 15 ถึง 30 & 10 ถึง 25 & 70 ถึง 120 \\
    อ้อยโรงงาน & $6.0 - 7.5$ & 24 ถึง 32 & 25 ถึง 35 & 55 ถึง 75 & 25 ถึง 45 & 20 ถึง 35 & 90 ถึง 150 \\
    \midrule
    \multicolumn{8}{l}{\textbf{3. กลุ่มไม้ผล (Fruit Trees / Orchards)}} \\
    \textbf{ทุเรียน (หมอนทอง/ชะนี)} & $\mathbf{5.5 - 6.5}$ & \textbf{25 ถึง 30} & \textbf{25 ถึง 32} & \textbf{50 ถึง 65} & \textbf{25 ถึง 45} & \textbf{25 ถึง 50} & \textbf{100 ถึง 180} \\
    มังคุด & $5.5 - 6.5$ & 24 ถึง 30 & 25 ถึง 32 & 55 ถึง 70 & 20 ถึง 35 & 15 ถึง 30 & 80 ถึง 140 \\
    มะม่วง (น้ำดอกไม้) & $5.5 - 7.0$ & 24 ถึง 32 & 24 ถึง 35 & 40 ถึง 60 & 20 ถึง 40 & 15 ถึง 30 & 90 ถึง 160 \\
    \bottomrule
\end{tabularx}%
}
\end{table}

\begin{table}[H]
\centering
\caption{เกณฑ์ธาตุอาหารรอง Ca Mg S จุลธาตุ และอาการผิดปกติของพืชเศรษฐกิจ}
\label{tab:crop_soil_secondary_micro}
\resizebox{\textwidth}{!}{%
\begin{tabularx}{1.18\textwidth}{p{2.6cm}cccp{3.2cm}Y}
    \toprule
    \textbf{ชนิดพืช} & \textbf{แคลเซียม Ca (mg/kg)} & \textbf{แมกนีเซียม Mg (mg/kg)} & \textbf{กำมะถัน S (mg/kg)} & \textbf{จุลธาตุสำคัญ (mg/kg)} & \textbf{สัญญาณเตือนภาวะผิดปกติและคำแนะนำ} \\
    \midrule
    ข้าวหอมมะลิ & 300 ถึง 600 & 80 ถึง 150 & 15 ถึง 30 & Zn: 1.5 ถึง 3.0, Fe: 50 ถึง 150 & หากดินน้ำขังนานและขาดสังกะสี ต้นข้าวจะแคระแกร็น ยอดสั้น \\
    ข้าวเจ้าชลประทาน & 350 ถึง 700 & 90 ถึง 180 & 20 ถึง 40 & Zn: 2.0 ถึง 4.0, Mn: 20 ถึง 60 & ระวังความเป็นพิษของเหล็กในดินกรดจัดแต่น้ำขัง \\
    ข้าวโพดเลี้ยงสัตว์ & 400 ถึง 800 & 100 ถึง 200 & 20 ถึง 35 & Zn: 2.0 ถึง 5.0, B: 0.5 ถึง 1.5 & ขาดสังกะสีเกิดแถบขาวขนานเส้นใบ ขาดโบรอนเมล็ดฟีบไม่เต็มฝัก \\
    มันสำปะหลัง & 200 ถึง 500 & 50 ถึง 120 & 10 ถึง 25 & Fe: 15 ถึง 40, Mn: 15 ถึง 50 & ขาดโพแทสเซียมหัวไม่ขยาย ขาดแมกนีเซียมใบแก่เหลืองระหว่างเส้นใบ \\
    อ้อยโรงงาน & 500 ถึง 1000 & 120 ถึง 250 & 25 ถึง 45 & B: 1.0 ถึง 2.0, Cu: 0.5 ถึง 1.5 & แคลเซียมช่วยเพิ่มความหวานและน้ำหนักลำ ป้องกันไส้กลวง \\
    \textbf{ทุเรียน} & \textbf{600 ถึง 1200} & \textbf{150 ถึง 300} & \textbf{25 ถึง 50} & \textbf{B: 1.0 ถึง 2.5, Zn: 2.0 ถึง 5.0} & \textbf{ขาดโบรอนหนามผลจีบผลเบี้ยว ขาดแคลเซียมเปลือกแตกเต่าเผา} \\
    มังคุด & 500 ถึง 900 & 120 ถึง 220 & 20 ถึง 40 & Fe: 20 ถึง 50, Zn: 1.5 ถึง 4.0 & แคลเซียมไม่พอเกิดอาการเนื้อแก้วและยางไหลซึมในเนื้อผล \\
    มะม่วง & 500 ถึง 1000 & 100 ถึง 200 & 20 ถึง 40 & B: 0.8 ถึง 2.0, Zn: 1.5 ถึง 4.0 & ขาดโบรอนและแคลเซียมส่งผลให้เนื้อผลด้านและผลแตกง่าย \\
    \bottomrule
\end{tabularx}%
}
\end{table}

\subsection{กลยุทธ์การตรวจวัดและจัดการสภาพแวดล้อมตามระยะการเจริญเติบโต}

การจัดการดินอย่างแม่นยำต้องสอดคล้องกับระยะพัฒนาการของพืช โดยเฉพาะไม้ผลมูลค่าสูงเช่นทุเรียน
\begin{enumerate}
    \item \textbf{ระยะฟื้นฟูต้นและแตกใบอ่อน (Flushing Stage)} ต้องการดินที่มีค่า pH 6.0 ถึง 6.5 ไนโตรเจนสูง (35 ถึง 45 มก./กก.) ร่วมกับแมกนีเซียมและสังกะสี เพื่อกระตุ้นการสังเคราะห์คลอโรฟิลล์และการขยายขนาดใบ
    \item \textbf{ระยะสะสมอาหารและกักโศกทำดอก (Flower Bud Differentiation)} ต้องควบคุมความชื้นดินให้ลดลงสู่ร้อยละ 25 ถึง 35 VWC ต่อเนื่อง 10 ถึง 14 วัน เพื่อกระตุ้นฮอร์โมนเปิดตาดอก พร้อมควบคุมไนโตรเจนให้อยู่ในระดับต่ำและยกระดับฟอสฟอรัสและโพแทสเซียม
    \item \textbf{ระยะบานดอกและผสมเกสร (Anthesis \& Fruit Set)} คุมความชื้นดินปานกลางร้อยละ 45 ถึง 55 VWC หลีกเลี่ยงการให้น้ำกระชาก เสริมธาตุโบรอนและแคลเซียมเพื่อส่งเสริมการงอกของหลอดละอองเรณู
    \item \textbf{ระยะขยายขนาดผลและสะสมเนื้อ (Fruit Swelling \& Bulking)} ความต้องการน้ำและโพแทสเซียมพุ่งสูงสุด (ความชื้นดินร้อยละ 55 ถึง 65 VWC และ K มากกว่า 120 มก./กก.) พร้อมรักษาค่า pH ไม่ให้ต่ำกว่า 5.5 เพื่อป้องกันการชะงักการดูดซึมแคลเซียม
\end{enumerate}

\section{สถาปัตยกรรมการประมวลผลสองขั้นตอนแบบ Edge Two-Stage Decoupling}

เพื่อแก้ปัญหาความคลาดเคลื่อนสะสมจากการเปลี่ยนแปลงอุณหภูมิและความชื้นดินตามรอบวัน ระบบ \textbf{SoilpH-HT} ได้รับการออกแบบสถาปัตยกรรมแบบสองขั้นตอน (Two-Stage Decoupling) ทำงานบนสมาร์ตโฟนโดยตรง

\begin{figure}[H]
\centering
\begin{tikzpicture}[
    node distance=1.6cm,
    block/.style={rectangle, rounded corners=2mm, draw=rbruNavy, fill=rbruNavy!8, thick, minimum width=5.0cm, minimum height=0.9cm, align=center},
    arrow/.style={-{Stealth[scale=1.2]}, thick, draw=rbruBlue}
]
    \node[block, fill=cyan!15] (sensor) {\textbf{หัววัด Soil Parameter Tester}\\สัญญาณดิบ ($\text{pH}_{\text{raw}}, T_{\text{soil}}, M_{\text{soil}}$) ผ่าน USB-C OTG};
    \node[block, fill=rbruTeal!20, below=0.8cm of sensor] (stage1) {\textbf{ขั้นตอนที่ 1 ฟิสิกส์หลักการแรก}\\Nernst Slope $\Delta\text{pH}_T$ + Dry Contact $\Delta\text{pH}_M$};
    \node[block, fill=rbruGold!20, below=0.8cm of stage1] (stage2) {\textbf{ขั้นตอนที่ 2 โครงข่ายประสาทเทียม PINN}\\Non-linear Cross-Coupled Residual MLP};
    \node[block, fill=green!20, below=0.8cm of stage2] (ui) {\textbf{แดชบอร์ด SoilpH-HT}\\แสดงผลสด ชดเชยเรียลไทม์ บันทึกภาพ HUD และพิกัด};

    \draw[arrow] (sensor) -- (stage1);
    \draw[arrow] (stage1) -- (stage2);
    \draw[arrow] (stage2) -- (ui);
\end{tikzpicture}
\caption{ผังการทำงานของสถาปัตยกรรมการประมวลผลสองขั้นตอนแบบ Edge Two-Stage Decoupling สำหรับวิเคราะห์ค่ากรด-ด่างดิน}
\label{fig:two_stage_arch}
\end{figure}

\section{สรุปสาระสำคัญประจำบทที่ 1}

การทำความเข้าใจรากฐานฟิสิกส์เคมีของการตรวจวัดค่ากรด-ด่างดิน เป็นหัวใจสำคัญในการออกแบบระบบตรวจวัดอัจฉริยะ โดยสามารถสังเคราะห์สาระสำคัญได้ดังนี้
\begin{enumerate}
    \item \textbf{ค่าความเป็นกรด-ด่างเป็นตัวแปรควบคุมหลัก} ค่า pH ในดินควบคุมการละลายของสารอาหาร หากดินมีสภาพเป็นกรดรุนแรง (pH ต่ำกว่า 5.0) จะปลดปล่อยอะลูมิเนียมที่เป็นพิษและตรึงฟอสฟอรัสอย่างรุนแรง
    \item \textbf{ความคลาดเคลื่อนจากอุณหภูมิและความชื้น} สัญญาณศักย์ไฟฟ้าเคมีของหัววัดแปรผันตรงตามอุณหภูมิสัมบูรณ์ตามสมการเนิร์นสต์ และเกิดความต้านทานสัมผัสสูงเมื่อความชื้นในดินลดต่ำกว่าร้อยละ 25
    \item \textbf{สถาปัตยกรรม Two-Stage Decoupling} การผสานฟิสิกส์หลักการแรกเข้ากับโครงข่ายประสาทเทียม PINN บนสมาร์ตโฟน ช่วยให้สามารถแยกสัญญาณรบกวนของอุณหภูมิและความชื้นออกจากค่ากรด-ด่างที่แท้จริงได้อย่างแม่นยำ
\end{enumerate}

\section*{คำถามทบทวนท้ายบทที่ 1}
\addcontentsline{toc}{section}{คำถามทบทวนท้ายบทที่ 1}

\begin{enumerate}
    \item จงอธิบายบทบาทของกัมมันตภาพของไฮโดรเจนไอออน ($a_{\text{H}^+}$) ในการนิยามค่าความเป็นกรด-ด่างของดิน
    \item เพราะเหตุใดอุณหภูมิดินที่สูงขึ้นเป็น 40 องศาเซลเซียส จึงทำให้ความชันเนิร์นสเตียนเพิ่มขึ้น และส่งผลกระทบต่อค่า pH ที่หัววัดอ่านได้อย่างไร
    \item จงอธิบายกลไกที่ทำให้ดินแห้ง (ความชื้นต่ำกว่าร้อยละ 25) ส่งผลให้เกิดความคลาดเคลื่อนในการวัดค่า pH
    \item จงเปรียบเทียบปริมาณปูนโดโลไมต์ที่ต้องใช้ในการยกระดับค่า pH ขึ้น 1.0 หน่วย ระหว่างดินทรายกับดินเหนียว พร้อมระบุเหตุผลทางปฐพีวิทยา
\end{enumerate}
